% Curriculum Vitae — Rajarshi SinhaRoy
% Build: pdflatex Rajarshi_SinhaRoy_CV.tex   (needs fontawesome5 + tex-gyre; both on Overleaf)
% signature.jpg is kept locally only (gitignored).
\documentclass[10pt,a4paper]{article}

\usepackage[T1]{fontenc}
\usepackage[utf8]{inputenc}
\usepackage{tgheros}
\renewcommand{\familydefault}{\sfdefault}
\usepackage[margin=0.6in]{geometry}
\usepackage{xcolor}
\usepackage{graphicx}
\usepackage{fontawesome5}
\usepackage{titlesec}
\usepackage{enumitem}
\usepackage[hidelinks]{hyperref}
\urlstyle{same}

\definecolor{rulegray}{gray}{0.55}
\definecolor{doigray}{gray}{0.45}

\pagestyle{empty}
\setlength{\parindent}{0pt}
\setlength{\parskip}{0pt}

\titleformat{\section}{\Large\scshape}{}{0pt}{}[{\color{rulegray}\titlerule[0.4pt]}]
\titlespacing*{\section}{0pt}{8pt}{19pt}

\setlist[itemize]{leftmargin=1.2em,label={\scriptsize\textbullet},itemsep=0pt,parsep=0pt,topsep=4pt}

\newcommand{\doi}[1]{\href{https://doi.org/#1}{\textcolor{doigray}{doi:\nolinkurl{#1}}}}
\newcommand{\sep}{\quad\textbar\quad}
\newcommand{\entry}[3]{\textbf{#1}\\\textit{#2}\hfill\textit{#3}\par\vspace{8pt}}
\newcommand{\pubhead}[1]{\textbf{#1}\par}

\begin{document}

% ---------- Header ----------
\begin{center}
  {\LARGE\bfseries Rajarshi SinhaRoy}\par\vspace{5pt}
  {\small
  \faMapMarker*\ Leipzig, Germany\sep
  \faEnvelope\ \href{mailto:sinharoy.rajarshi@uni-leipzig.de}{sinharoy.rajarshi@uni-leipzig.de}\par\vspace{3pt}
  \faLinkedin\ \href{https://www.linkedin.com/in/rajarshi-sinharoy}{LinkedIn}\sep
  \faGithub\ \href{https://github.com/Zadak-H}{GitHub}\sep
  \faGraduationCap\ \href{https://scholar.google.com/citations?user=CQJhB0IAAAAJ}{Google Scholar}}
\end{center}

% ---------- Research Profile ----------
\section{Research Profile}
{\fontsize{10.95}{13.8}\selectfont
PhD researcher in Informatics at the University of Leipzig (Germany), specializing in Artificial Intelligence
for Drug Discovery. My research integrates deep learning, bioinformatics, and structural biology to design
interpretable models that predict enzyme function, protein stability, and structure–function relationships.
With expertise in geometric deep learning, molecular simulation, and biomolecular modeling, I aim to bridge
algorithmic innovation and molecular biophysics for predictive, explainable discovery of therapeutic enzymes
and drug targets.\par}

% ---------- Education ----------
\section{Education}
{\large\bfseries PhD in Informatics}\\
\textit{Institute of Drug Discovery, University of Leipzig, Germany}\hfill\textit{Nov 2023 – Present}\par
\textbf{Focus:} AI for mRNA delivery, Geometric Deep Learning, Protein Modeling, Brownian Dynamics, Molecular Dynamics
and Computational Biology.\par
\textbf{Supervisors:} Prof.\ Dr.\ Jens Meiler; Dr.\ Georg Künze.\par
\vspace{9pt}
{\large\bfseries M.Sc.\ in Computer Science}\\
\textit{St.\ Xavier’s College (Autonomous), Kolkata, India}\hfill\textit{2020 – 2022}\par
CGPA 8.36 \enspace\textbullet\enspace Dissertation: \textit{Optimization in Image Processing and Deep Learning.}\par
\vspace{9pt}
{\large\bfseries B.Sc.\ in Computer Science}\\
\textit{Government General Degree College, Singur, India}\hfill\textit{2017 – 2020}\par
CGPA 8.83 \enspace\textbullet\enspace Focus: Artificial Intelligence, Database Systems, Applied Computing.\par

% ---------- Experience ----------
\section{Research \& Professional Experience}
\entry{PhD Researcher — Artificial Intelligence for Drug Discovery}{University of Leipzig, Germany}{Nov 2023 – Present}
\begin{itemize}
  \item Collaborative project with \textbf{BioNTech} on enhancing mRNA delivery therapeutics through computational design.
  \item Developing \textbf{SE(3)-equivariant GNNs} for enzyme pH-optima prediction.
  \item Designing \textbf{MPNN–Transformer} architectures for pH-dependent protein sequence generation.
  \item Conducting molecular-dynamics simulations and bioinformatics analyses for interpretable AI models.
\end{itemize}
\vspace{2pt}
\entry{Academic Researcher — Bioinformatics and AI Lab}{Kolkata, India}{2021 – 2023}
\begin{itemize}
  \item Designed hybrid \textbf{GAN + DCNN} architectures for MRI-based Alzheimer’s progression prediction.
  \item Built preprocessing, augmentation, and evaluation pipelines ensuring reproducibility.
\end{itemize}
\vspace{2pt}
\entry{Teaching Assistant — Computer Science Department}{Techno Main Salt Lake, India}{Jan 2023 – Oct 2023}
\begin{itemize}
  \item Assisted undergraduate courses in Computer Science and Data Analytics.
  \item Conducted lab tutorials and supervised ML/AI-based student projects.
\end{itemize}
\vspace{2pt}
\entry{Developer Intern — EasyWay Technologies}{Kolkata, India}{Jan 2022 – Apr 2022}
\begin{itemize}
  \item Developed and deployed backend services (PHP, MySQL) and data visualization modules for a client-facing
        analytics platform.
  \item Implemented analytics automation and optimized API performance.
\end{itemize}
\vspace{2pt}
\entry{Software Developer Intern — Unique Start}{Kolkata, India}{May 2021 – Oct 2021}
\begin{itemize}
  \item Developed REST APIs and front-end integrations for a scalable data platform.
  \item Improved scalability by optimizing SQL queries and adding a data caching layer, reducing API latency.
\end{itemize}

% ---------- Publications ----------
\section{Publications \& Conferences}
\pubhead{Peer-Reviewed Publications}
\begin{itemize}
  \item R. SinhaRoy, C. Clauss, I. Ivanikov, and G. Künze, “Predicting Enzyme pH Optima from Structure Using Equivariant
        Graph Neural Networks,” \textit{Journal of Chemical Information and Modeling}, 66(16), 9858–9873, 2026.
        \doi{10.1021/acs.jcim.6c01482}.
  \item R. SinhaRoy and A. Sen, “A Hybrid Deep Learning Framework to Predict Alzheimer’s Disease Progression Using
        GAN and DCNN,” \textit{Arabian Journal for Science and Engineering}, 2023. \doi{10.1007/s13369-023-07973-9}.
  \item R. Sinharoy and A. Sen, “Cardiovascular Disease Prediction Using Ensemble Classification Algorithm in Machine
        Learning,” \textit{International Journal of Soft Computing}, 2022. \doi{10.21917/ijsc.2022.0366}.
  \item R. Sinharoy and S. Sarkhel, “Air Quality Index Prediction in Realtime Using SVM,” \textit{International Journal of Innovative
        Research in Physics}, 2021. \doi{10.15864/ijiip.3106}.
  \item R. SinhaRoy, “A Study on the Journey of NLP Models: From Symbolic NLP to BERT,” \textit{IJSRCSEIT}, 2021.
        \doi{10.32628/CSEIT217688}.
\end{itemize}
\vspace{4pt}
\pubhead{Conference Presentations}
\begin{itemize}
  \item R. SinhaRoy, “pHoptNN: Predicting Enzyme \textit{pH}-Optima from Structure Using Equivariant Graph Neural Networks,”
        Oral Presentation, \textit{EUROSETTACON 2025}.
  \item S. Sen, S. Sonali, R. SinhaRoy et al., “Overlapped Fingerprint Separation Using Graph-Based Model,” in \textit{IEEE
        Silchar Subsection Conference (SILCON)}, 2022. \doi{10.1109/SILCON55242.2022.10028828}.
\end{itemize}

% ---------- Skills ----------
\section{Technical Skills}
{\linespread{1.18}\selectfont
\textbf{Programming:} Python • C/C++ • Java • SQL • MATLAB • PHP • HTML/CSS/JS • Git\par
\textbf{ML/DL:} TensorFlow • Keras • PyTorch • GNNs • CNNs • GANs • MPNNs • Diffusion Models\par
\textbf{Bio-Software:} RFDiffusion • ProteinMPNN • GROMACS • AMBER • VMD • AlphaFold • PyMOL\par
\textbf{Research Areas:} Computational Biology • Bioinformatics • Data Mining • Molecular Dynamics\par}

% ---------- Signature ----------
\vfill
\begin{minipage}[b]{0.5\textwidth}
October 30, 2026\\
Leipzig, DE
\end{minipage}%
\hfill
\begin{minipage}[b]{0.25\textwidth}
\raggedleft
\IfFileExists{signature.jpg}{\includegraphics[width=\linewidth]{signature.jpg}\\[-4pt]}{\vspace{1.2cm}}
\rule{\linewidth}{0.4pt}\\
Rajarshi SinhaRoy
\end{minipage}
\vspace{6.15cm}

\end{document}
